本文针对北京工业大学耿丹学院本科生毕业设计论文的排版需求，开发了一套基于 \LaTeX{} 的自动化模板系统。主要工作与结论如下：

（1）系统梳理了学校毕业设计 Word 模板的全部格式规范，建立了完整的排版参数清单，涵盖页面设置、字体字号、各级标题、正文段落、图表公式、参考文献等各个方面。

（2）基于 ctex、titlesec、gbt7714 等成熟宏包，设计并实现了 \verb|gengdan-thesis.cls| 文档类文件，将复杂的格式配置封装为开箱即用的模板，大幅降低了学生的排版门槛。

（3）通过多平台、多场景的实验测试，验证了模板在排版精度、编译效率、可移植性等方面的优良性能，证明其完全满足学校毕业设计论文的提交要求。

未来可在以下方面进一步改进：增加对插图批量处理的辅助脚本；开发基于 VS Code 或 TeXstudio 的一键编译配置；针对不同专业（如艺术、经管类）的特殊格式需求提供扩展分支。